% !TeX root = ../tesis.tex
%%%%%%%%%%%%%%%%%%%%%%%%%%%%%%%%%%%%%%%%%%%%%%%%%%%%%%%%%%%%%%%%%%%
% SECCIÓN 5.1 — LA RED METROPOLITANA COMO OBJETO DE ANÁLISIS
% Issue #24 [C-7]
%%%%%%%%%%%%%%%%%%%%%%%%%%%%%%%%%%%%%%%%%%%%%%%%%%%%%%%%%%%%%%%%%%%
\section{La Red Metropolitana como Objeto de Análisis}
\label{sec:red-objeto-analisis}

El diagnóstico del sistema de transporte público de la \zmvm{} enfrenta un
problema de escala: la red integra once modalidades de transporte masivo
distribuidas sobre un territorio metropolitano que trasciende la \cdmx{} e
incorpora municipios del Estado de México en una continuidad urbana funcional
\citep{semovi2020diagnostico}. Los indicadores de movilidad son difíciles de
aplicar de forma homogénea a cada sistema de transporte, y este problema no
radica en la falta de datos: el Gobierno de la \cdmx{} libera sistemáticamente
la información de sus redes a través del Portal de Datos Abiertos
\citep{adip_gtfs_2024}. La barrera real es la heterogeneidad estructural de
esos datos: once sistemas con esquemas, frecuencias, geometrías y atributos
distintos que, sin una herramienta capaz de integrarlos bajo un modelo común,
no pueden analizarse como red sino solo como inventario.

Este es el problema que el \textbf{Apímetro} y el \textbf{Modelo VFT},
desarrollados y documentados en el Capítulo~4, están diseñados para resolver
\citep{cabrera2025apimetro, cabrera2026VFTModel}. Al representar la red como
un grafo dirigido y ponderado —donde cada estación es un nodo y cada segmento
operativo una arista con peso temporal— los once sistemas quedan integrados
bajo un modelo matemático común \citep{newman2010networks}. Sobre esa
representación es posible calcular propiedades estructurales independientes de
la demanda: la cobertura territorial de la red, la concentración de flujos en
los nodos de transferencia, la eficiencia de los trayectos y la fragilidad de
los puntos críticos.

El presente capítulo aplica esa herramienta al sistema de transporte actual de
la \zmvm{} para responder tres preguntas estructurales: ¿qué territorio cubre
la red y dónde se concentran los flujos? ¿Cuánto tarda en promedio
desplazarse entre cualquier par de estaciones y qué nodos articulan la mayor
parte de esos trayectos? ¿Qué tan frágil es la red ante la pérdida de sus
nodos más críticos? Las respuestas se organizan en tres fases de análisis
—indicadores topológicos, ponderación dinámica y análisis de red— y culminan
en un diagnóstico territorial que identifica las zonas de mayor déficit
estructural como punto de partida para el corredor anillar hipotético del
Capítulo~6.

El Cuadro~\ref{tab:resumen-indicadores-cap5} sintetiza los ocho indicadores
que componen el análisis, su fase de cálculo y su estado de avance.

\begin{table}[H]
    \centering
    \caption[Indicadores del Modelo VFT por fase de cálculo]{Resumen de los
    indicadores del \textbf{Modelo VFT}, su fase de cálculo, dependencias y
    estado de avance en esta investigación.}
    \label{tab:resumen-indicadores-cap5}
    \begin{tabularx}{\textwidth}{|l|c|X|l|}
        \hline
        \textbf{Indicador} & \textbf{Fase} & \textbf{Dependencias} &
        \textbf{Estado} \\
        \hline
        Validación de Impedancia      & 1 & Grafo base + fallbacks       & Calculado \\
        Cobertura Espacial ($C$)      & 1 & Geometría de estaciones      & Calculado \\
        Fuerza Capilar ($FC$)         & 1 & Grafo base                   & Calculado \\
        Factor de Desviación ($DI$)   & 1 & Grafo base + impedancia      & Calculado \\
        \hline
        Coef. de Fricción Vial ($CF$)     & 2 & Fase 1 completa          & Calculado \\
        Penalización por Transf. ($T_b$)  & 2 & Fase 1 + $CF$            & Calculado \\
        Node Score ($\text{VFT\_score}$)  & 2 & Fase 1 + datos MV1/MV2   & En diseño \\
        \hline
        Tiempo de Viaje Prom. ($T$)   & 3 & Fases 1 y 2 completas        & Calculado \\
        Centralidad Interm. ($B$)     & 3 & Fases 1 y 2 + $T$            & Calculado \\
        Vulnerabilidad ($\Delta E$)   & 3 & Fases 1, 2 y $B$             & Calculado \\
        \hline
    \end{tabularx}
    \fuentefigura{Elaboración propia con base en \texttt{VFTModel}
    \citep{cabrera2026VFTModel}.}
\end{table}
